% !TEX program = xelatex
\documentclass[11pt]{article}
\usepackage{fontspec}
\usepackage{amsmath,graphicx}
\usepackage[backend=biber]{biblatex}
\newcommand{\R}{\mathbb{R}}

\title{Edge cases for protection}
\author{A. Writer}

\begin{document}
\maketitle

\section{Introduction}\label{sec:intro}
This paragraph cites \cite{knuth84,lamport94} and refers to Section~\ref{sec:intro} and Eq.~\eqref{eq:one}.
Inline maths $a^2 + b^2 = c^2$ and \(x \in \R\) must stay as they are. A percent sign \% is not a comment, % but this is
and so is the rest of this line.

\begin{equation}\label{eq:one}
  E = mc^2
\end{equation}

\[ \sum_{i=1}^{n} i = \frac{n(n+1)}{2} \]

We use \textbf{bold words} and \emph{emphasised words} inside prose, plus \texttt{code\_name} and \verb|raw|text|.
A footnote\footnote{Footnotes are prose too.} and a URL \url{https://example.com/a_b}.

\begin{verbatim}
verbatim text is never changed
\end{verbatim}

\begin{figure}[ht]
  \centering
  \includegraphics[width=0.8\linewidth]{figures/plot.pdf}
  \caption{A caption with words that can be rewritten.}
  \label{fig:plot}
\end{figure}

\begin{tabular}{l|r}
  Name & Value \\ \hline
  alpha & 1 \\
\end{tabular}

$$ \int_0^\infty e^{-x}\,dx = 1 $$

Special characters: ~ & ^ _ { } and a nested \textbf{\emph{\textsc{group}}}.

\printbibliography
\end{document}
